\section{模型评价与推广}

\subsection{模型优点}

（1）模型体系完整，覆盖了从基础功率平衡到多场景优化调度的完整分析链条。问题一至问题四形成了"满负荷基准分析$\rightarrow$离散调节优化$\rightarrow$连续调节优化$\rightarrow$离网储能配置"的递进式研究框架，每一步都在前一步基础上放松约束或增加决策维度，逻辑清晰、层次分明。

（2）优化模型具有严格的数学基础和高效的求解性能。问题二采用0-1整数规划（MILP），问题三采用线性规划（LP），均为凸优化问题，可保证全局最优解的获取。采用PuLP/CBC求解器，24个时段的调度问题在秒级时间内完成求解，具备实际工程应用的计算效率。

（3）多场景分析框架具有较强的实用价值。通过6种风电$\times$4种光伏共24种场景的组合分析，全面覆盖了园区可能面临的风光出力条件，为园区全年运行策略制定提供了系统性的决策支持。每种场景代表15天的假设使得全年统计分析具有明确的物理意义。

（4）经济性分析维度丰富，为工程决策提供了定量依据。模型不仅计算了吨氨成本，还量化了电网支撑价值（13.7\%）、储能边际效益、连续调节相对离散调节的成本节约（6.3\%）等关键经济指标，为园区投资决策和运行模式选择提供了科学依据。

\subsection{模型不足}

（1）时间分辨率限制。模型以1小时为最小调度单元，无法捕捉分钟级的风光功率波动和设备启停过渡过程。实际运行中，电解槽的冷启动时间、功率爬坡速率等动态特性可能影响调度方案的可行性。

（2）设备效率简化。模型假设电解槽在10\%--100\%负荷范围内效率恒定（线性产氢），实际上电解槽在低负荷运行时效率会下降。引入非线性效率曲线将使问题变为非线性规划，增加求解难度但可提高结果精度。

（3）不确定性处理不足。模型采用确定性场景分析，24种场景虽然覆盖了不同风光水平，但未考虑场景间的转移概率和日内风光预测误差。引入随机规划或鲁棒优化方法可提高调度方案对不确定性的适应能力。

\subsection{模型推广}

本文建立的优化调度模型框架具有良好的可扩展性，可推广至以下场景：（1）多园区协同调度——多个绿电直连园区通过区域电网互联，实现容量互济和负荷互补；（2）氢储能耦合——将制氢产物作为长时储能介质，在风光过剩时制氢储存、在风光不足时燃料电池发电；（3）碳交易市场耦合——将绿氨的碳减排效益纳入目标函数，在碳价较高时增加产量以获取碳交易收益。这些扩展方向将进一步提升模型的工程应用价值和学术研究深度。
